\section{Related Work}
\label{sec:relwork}

\paragraph{Proof rules for almost-sure termination.}
The classical approach to AST verifies a \emph{variant function} that decreases with a uniformly positive probability, going back to McIver and Morgan~\cite{DBLP:series/mcs/McIverM05}; their rule is complete only when the reachable state space is finite. Supermartingale-based rules relax this requirement. McIver et al.~\cite{10.1145/3158121} give a rule in which a variant, itself required to be a supermartingale, decreases with a state-dependent probability and satisfies an additional progress condition; this rule also handles demonic nondeterminism and extends the class of provably AST programs beyond what bounded-variant arguments allow. Majumdar and Sathiyanarayana~\cite{DBLP:journals/pacmpl/MajumdarS25} recently gave the first proof rule for AST that is both sound and relatively complete, formulated over probabilistic control-flow graphs (pCFGs). They further show that the earlier rule of McIver et al.~\cite{10.1145/3158121} enjoys the same relative-completeness guarantee when read at the program level. Our rule is derived from the pCFG rule of~\cite{DBLP:journals/pacmpl/MajumdarS25} and inherits its soundness and relative completeness, but it is stated directly over distribution-transformer semantics and factors the supermartingale and variant obligations into two independent certificates rather than a single combined one. This factorization is what lets rejection loops reset the variant without disturbing the supermartingale argument, which is central to our Fast Loaded Dice Roller (FLDR) case study. Qualitative AST is the property we target throughout; positive AST, which additionally bounds the expected runtime, is a strictly stronger and differently characterized property that we do not address.

\paragraph{Operational and denotational models.}
pCFGs and Markov decision processes are the standard operational substrate for termination and reachability analysis of probabilistic and randomized systems~\cite{DBLP:books/wi/Puterman94,DBLP:books/sp/BaierK08}. Tools such as Storm~\cite{DBLP:journals/sttt/HenselJKQV22} and refined analyses of visiting times and stationary behavior~\cite{DBLP:conf/tacas/MertensKQW24,DBLP:journals/jcss/JungesK0W21} operate at this level. Viewing probabilistic systems as \emph{distribution transformers} rather than as single-state transition systems has recently been fruitful for both Markov chains and MDPs: Aghamov et al.~\cite{DBLP:journals/corr/abs-2406-15087,DBLP:conf/birthday/AghamovBKNOPV25} model-check Markov chains through their action on distributions, and Akshay et al.~\cite{DBLP:conf/cav/AkshayCMZ23,DBLP:conf/ijcai/0001CMZ24} synthesize invariants and certified policies for MDPs under distributional reach-avoidance objectives; distribution-based objectives for MDPs are also studied in~\cite{DBLP:conf/lics/AkshayGV18}. Our rule applies the same distribution-transformer perspective one level higher, directly to program syntax rather than to an underlying automaton, which avoids the location-indexed bookkeeping that a translation to pCFGs or MDPs otherwise imposes.

\paragraph{Denotational semantics and program logics for probabilistic programs.}
The denotational, post-distribution-transformer semantics we build on originates with Kozen~\cite{DBLP:conf/focs/Kozen79,DBLP:conf/stoc/Kozen83} and underlies the weakest-precondition and weakest pre-expectation calculi of McIver and Morgan~\cite{DBLP:series/mcs/McIverM05} and Kaminski~\cite{DBLP:phd/dnb/Kaminski19}, in the tradition of Dijkstra's predicate-transformer semantics~\cite{DBLP:journals/cacm/Dijkstra75,DBLP:books/daglib/0067387}. Hoare-style program logics for probabilistic programs were introduced by den Hartog and de Vink~\cite{DBLP:journals/ijfcs/HartogV02} and extended by the Ellora framework~\cite{DBLP:conf/esop/BartheEGGHS18}; more recent calculi include quantitative strongest-postcondition reasoning~\cite{DBLP:journals/pacmpl/ZhangK22}, its hyperproperty generalization~\cite{DBLP:journals/corr/abs-2404-05097}, and modal separation logic for conditional probability~\cite{DBLP:journals/pacmpl/Li0H23}. These logics target partial correctness or information flow rather than AST; our rule complements them by supplying the missing termination argument needed to lift a partial-correctness proof to total correctness.

\paragraph{Verifying exact sampling algorithms.}
Exact discrete samplers built from fair coin flips, following the classical constructions of von Neumann~\cite{vonNeumann1951} and Lumbroso's Fast Dice Roller (FDR)~\cite{DBLP:journals/corr/abs-1304-1916}, have received growing verification attention. Saad et al.\ introduce the Fast Loaded Dice Roller (FLDR) as a near-optimal sampler for arbitrary discrete distributions~\cite{DBLP:conf/aistats/SaadFRM20}; the original paper gives an informal termination argument, while our contribution supplies a proof in the program logic. Bagnall, Stewart, and Banerjee's Zar compiler~\cite{DBLP:journals/pacmpl/Bagnall0023,10.1145/3591220} formally verifies samplers compiled from probabilistic programs with loops and conditioning in Coq, targeting the random-bit model; generating-function techniques have likewise been used to reason about probabilistic programs and about samplers such as the fast dice roller~\cite{DBLP:conf/lopstr/KlinkenbergBKKM20,DBLP:conf/cav/ChenKKW22,haase2026generatingfunctionsmeetoccupation}. Exact and guaranteed-bound inference engines such as PSI~\cite{DBLP:conf/cav/GehrMV16,DBLP:conf/pldi/GehrSV20}, the approach of Beutner et al.~\cite{DBLP:conf/pldi/BeutnerOZ22}, and scalable exact inference for discrete programs~\cite{DBLP:journals/pacmpl/HoltzenBM20} address the complementary problem of computing or bounding the output distribution rather than proving termination. Closest to our work is Zilken and Batz et al.'s framework of distributional invariants~\cite{zilken2026verifyingsamplingalgorithmsdistributional}, which proves partial correctness of FDR-style samplers by tracking distributional invariants under the same post-distribution-transformer semantics we use. That work establishes that FDR's output is uniformly distributed \emph{conditioned on termination}; our AST proofs for FDR and FLDR are the missing ingredient for total correctness, and our rule is deliberately phrased so that its distributional invariant obligation can reuse the one already constructed for partial correctness.

\paragraph{Nondeterminism.}
Several lines of work verify probabilistic programs or Markov models in the presence of demonic or angelic nondeterminism~\cite{10.1145/3158121,DBLP:journals/pacmpl/BatzBKW24,Wang2018ADS,Zilberstein2025,10.1145/3743131,Zilberstein_2025}. Our proof rule targets probabilistic while-loops without nondeterministic choice; extending the distribution-transformer formulation to resolve nondeterminism, e.g., by adapting strategy-based semantics as in~\cite{DBLP:journals/pacmpl/BatzBKW24}, is an interesting direction we do not pursue here.
